\documentclass[11pt]{article}
\usepackage[margin=2.2cm]{geometry}
\usepackage{amsmath,amssymb,booktabs,enumitem}
\usepackage[hidelinks]{hyperref}
\setlist{nosep,leftmargin=*}
\setlength{\parindent}{0pt}

\newcommand{\dt}{\tilde\delta}
\newcommand{\nb}{\bar n}
\newcommand{\R}{\mathbb R}

\title{1 --- Departure from CSR: $\dt$ and null tables}
\date{}

\begin{document}
\maketitle

%==============================================================================
\section{Theory}

\subsection*{Setting}
\begin{itemize}
  \item Window $W=[0,1]^2$
  \item $\nb$ = expected number of points in $W$.
  \item Point cloud $x = \{x_1,\dots,x_n\} \subset W$; $n(x) = |x|$ denotes realised number of points
\end{itemize}

\subsection*{Estimator}
\begin{itemize}
  \item Centred $L$-function of a cloud $x$:
  \[
    \hat\phi_x(r) = \sqrt{\hat K_x(r)/\pi} - r,
    \qquad
    (= 0 \text{ under CSR})
  \]
  where $\hat K_x(r)$ is Ripley's edge-corrected $K$ function evaluated at radius $r$.
  \item Evaluated on a fixed grid $0<r_1<\dots<r_K=0.25$
\end{itemize}

\subsection*{Radius set}
\begin{itemize}
    \item For a cloud with $n = n(x)$ points we evaluate $\hat\phi_x(r)$ on
        \[
        R(n) = \{r_k : r_k \ge r_{\min}(n)\},
        \]
    where $r_{\min}(n)$ is the smallest radius with at least $p$ expected pairs closer than $r$ under CSR ($\mathbb P_{0,n}$):
        \[
        r_{\min}(n) = \min\{r : \mu(r;n) \ge p\},
        \qquad
        \mu(r;n) = \tfrac12 n(n-1)\pi r^2 .
        \]
    \item Use use $p=5$. Then
        \[
        r_{\min}(n) = \sqrt{\frac{2p}{\pi n(n-1)}} \approx \frac{1.78}{n}
        \]
\end{itemize}

\subsection*{Null distribution}
We compute the null distribution at every $n(x)$ on a grid.
\begin{itemize}
  \item Generate $12\,000$ CSR clouds from $\mathbb P_{0,n}$
  \item Estimate null moments for $r \in R(n)$ using half of the CSR clouds:
  \[
    m_0(r;n) = \mathbb E_{0,n}\big[\hat\phi_X(r)\big],
    \qquad
    s_0(r;n) = \mathrm{sd}_{0,n}\big[\hat\phi_X(r)\big]
  \]
  \item Using remaining half of CSR clouds, compute maximum deviation of each cloud $x$:
  \[
    M_n(x) = \max_{r\in R(n)} \frac{|\hat\phi_x(r) - m_0(r;n)|}{s_0(r;n)}
  \]
  \item Determine critical value: $c_{95}(n)$ = 95\% quantile of $M_n$ over these clouds.
  \item Repeated for every $n(x)$ on grid, then smooth $m_0$, $s_0$, $c_{95}$ across $n(x)$ (cubic splines in $\log n$)
\end{itemize}

\subsection*{Studentised global envelope (Myllym\"aki et al.\ 2017)}
\begin{itemize}
  \item Reject CSR at 5\% iff $M_n(x) > c_{95}(n)$, $n = |x|$
  \item We define the "distance from CSR" $\dt$:
    \[
      \dt(\theta,\nb) = \frac{1}{c_{95}(\nb)} \max_{r\in R(\nb)} \frac{|L_\theta(r) - r|}{s_0(r;\nb)}
    \]
  \item $\dt = 0$ at CSR. $\dt = 1$ is on the rejection boundary
  \item This property is known before simulation, so is usable as a sweep coordinate
  \item At fixed shape, $\dt$ grows with $\nb$ ($s_0$ shrinks)
\end{itemize}

\end{document}
